% 康奈尔听课笔记模板（配色取自王道课件）
\documentclass[10.5pt,a4paper,oneside]{article}
\usepackage[a4paper,top=2.3cm,bottom=2.2cm,left=5.7cm,right=1.5cm,
  marginparwidth=4.25cm,marginparsep=0.55cm,headheight=14pt,footskip=1.1cm]{geometry}
\usepackage{xeCJK}
\usepackage{xeCJKfntef}
\xeCJKDeclareCharClass{FullLeft}{"201C,"2018}
\xeCJKDeclareCharClass{FullRight}{"201D,"2019}
\xeCJKDeclareCharClass{CJK}{"2460->"24FF,"2605,"2606}
\setmainfont{Noto Sans CJK SC}[BoldFont={Noto Sans CJK SC Bold},ItalicFont={Noto Sans CJK SC},ItalicFeatures={FakeSlant=0.18}]
\setsansfont{Noto Sans CJK SC}[BoldFont={Noto Sans CJK SC Bold}]
\setmonofont{DejaVu Sans Mono}[Scale=0.9]
\setCJKmainfont{Noto Sans CJK SC}[BoldFont={Noto Sans CJK SC Bold},ItalicFont={Noto Sans CJK SC},ItalicFeatures={FakeSlant=0.18}]
\setCJKsansfont{Noto Sans CJK SC}[BoldFont={Noto Sans CJK SC Bold}]
\setCJKmonofont{Noto Sans Mono CJK SC}[Scale=0.9]
\linespread{1.32}
\usepackage{amsmath,amssymb}
\usepackage[dvipsnames,table]{xcolor}
\usepackage{graphicx}
\usepackage{adjustbox}
\usepackage{tikz}
\usetikzlibrary{calc}
\usepackage[most]{tcolorbox}
\usepackage{eso-pic}
\usepackage{fancyhdr}
\usepackage{titlesec}
\usepackage{needspace}
\usepackage{enumitem}
\usepackage{xltabular,booktabs,array,colortbl}
\usepackage{marginnote}
\reversemarginpar
\renewcommand*{\raggedleftmarginnote}{}
\renewcommand*{\raggedrightmarginnote}{}
\usepackage{caption}
\usepackage{titletoc}
\usepackage{fancyvrb}
\usepackage{framed}
\usepackage[section]{placeins}
\usepackage{hyperref}
\hypersetup{colorlinks=true,linkcolor=wdblue,urlcolor=wdblue,pdftitle={1.1
课后习题：数据结构的基本概念},pdfauthor={听课笔记}}

% ---------- 课件配色 ----------
\definecolor{wdblue}{HTML}{4472C4}
\definecolor{wdlightblue}{HTML}{5B9BD5}
\definecolor{wdtitle}{HTML}{4A86C8}
\definecolor{wddark}{HTML}{1F3864}
\definecolor{wdorange}{HTML}{ED7D31}
\definecolor{wdgray}{HTML}{A5A5A5}
\definecolor{wdgold}{HTML}{FFC000}
\definecolor{wdgreen}{HTML}{70AD47}
\definecolor{wddeepgreen}{HTML}{3F7F2A}
\definecolor{wdred}{HTML}{C00000}
\definecolor{wdpurple}{HTML}{7030A0}
\definecolor{cuebg}{HTML}{F2F6FC}
\definecolor{codebg}{HTML}{F7F9FC}
\color[HTML]{262626}

% ---------- 行内样式 ----------
\newcommand{\hlred}[1]{\textcolor{wdred}{\bfseries #1}}
\newcommand{\hlblue}[1]{\textcolor{wdblue}{\bfseries #1}}
\newcommand{\hlgreen}[1]{\textcolor{wddeepgreen}{\bfseries #1}}
\newcommand{\hlorange}[1]{\textcolor{wdorange}{\bfseries #1}}
\newcommand{\hlpurple}[1]{\textcolor{wdpurple}{\bfseries #1}}
\newcommand{\uln}[1]{\CJKunderline[format=\color{wdred},thickness=0.9pt,depth=0.28em]{#1}}
\newcommand{\wavy}[1]{\CJKunderwave[format=\color{wdorange}]{#1}}
\newcommand{\hlmark}[1]{\CJKunderline[format=\color{wdgold!55},thickness=0.55em,depth=-0.05em]{#1}}
\newcommand{\tagbox}[2]{\tcbox[on line,colback=#1,colframe=#1,coltext=white,boxsep=0.5pt,left=3pt,right=3pt,top=0.5pt,bottom=0.5pt,arc=2pt,fontupper=\small\bfseries]{#2}}

% ---------- 康奈尔版式背景：左侧线索栏 ----------
\newlength{\cuestripw}\setlength{\cuestripw}{5.15cm}
\AddToShipoutPictureBG{%
  \begin{tikzpicture}[remember picture,overlay]
    \fill[cuebg] (current page.north west) rectangle ($(current page.south west)+(\cuestripw,0)$);
    \draw[wdblue,line width=1.2pt] ($(current page.north west)+(\cuestripw,-1.55cm)$) -- ($(current page.south west)+(\cuestripw,1.55cm)$);
    \node[anchor=west,text=wdblue,font=\small\bfseries] at ($(current page.north west)+(0.75cm,-1.25cm)$) {线索 · 问题};
    \node[anchor=west,text=wdblue,font=\small\bfseries] at ($(current page.north west)+(\cuestripw+0.55cm,-1.25cm)$) {笔记 · 老师主线 + 补充讲解};
  \end{tikzpicture}}

% ---------- 页眉页脚 ----------
\pagestyle{fancy}\fancyhf{}
\renewcommand{\headrulewidth}{0pt}
\fancyhead[R]{\small\color{wdgray}第1章 绪论}
\fancyfoot[C]{\small\color{wdgray}\thepage}
\fancyfoot[R]{\small\color{wdorange}1.1
课后习题：数据结构的基本概念 · 听课笔记}

\fancypagestyle{plain}{\fancyhf{}\fancyfoot[C]{\small\color{wdgray}\thepage}\fancyfoot[R]{\small\color{wdorange}1.1
课后习题：数据结构的基本概念 · 听课笔记}}

% ---------- 标题：仿课件的蓝色标题 + 六色条 ----------
\newcommand{\sixbar}{\par\nointerlineskip\vspace{3pt}\centering
  \begin{tikzpicture}\foreach \c [count=\i] in {wdlightblue,wdorange,wdgray,wdgold,wdblue,wdgreen}
     {\fill[\c] ({(\i-1)*1.05},0) rectangle ({(\i-1)*1.05+0.92},0.11);}\end{tikzpicture}\par}
\titleformat{\section}[block]{\centering\Large\bfseries\color{wdtitle}}{}{0pt}{}[\sixbar]
\titlespacing*{\section}{0pt}{16pt}{10pt}
\pretocmd{\section}{\Needspace{8\baselineskip}}{}{}
\titleformat{\subsection}{\Needspace{7\baselineskip}\large\bfseries\color{wddark}}{\textcolor{wdblue}{\rule[-0.15em]{4pt}{1.05em}}\hspace{6pt}}{0pt}{}
\titlespacing*{\subsection}{0pt}{12pt}{5pt}
\titleformat{\subsubsection}{\Needspace{5\baselineskip}\normalsize\bfseries\color{wdblue}}{\textcolor{wdorange}{\large\ensuremath{\blacktriangleright}}\hspace{4pt}}{0pt}{}
\titlespacing*{\subsubsection}{0pt}{8pt}{3pt}
\setcounter{secnumdepth}{0}
\setlength{\parindent}{0pt}\setlength{\parskip}{5pt}
\setlist{nosep,leftmargin=1.5em,itemsep=2pt}
\setlist[itemize,1]{label=\textcolor{wdblue}{\small\ensuremath{\bullet}}}
\setlist[itemize,2]{label=\textcolor{wdorange}{\scriptsize\ensuremath{\blacktriangleright}}}
\setlist[enumerate,1]{label=\textcolor{wdblue}{\bfseries\arabic*.}}
\providecommand{\tightlist}{\setlength{\itemsep}{2pt}\setlength{\parskip}{0pt}}

% ---------- 目录 ----------
\setcounter{tocdepth}{3}
\titlecontents{section}[0pt]{\addvspace{5pt}\bfseries\color{wddark}}{}{}{\titlerule*[0.7pc]{\color{wdgray!60}.}\contentspage}
\titlecontents{subsection}[1.6em]{\addvspace{1pt}\small\color{black!75}}{}{}{\titlerule*[0.7pc]{\color{wdgray!40}.}\contentspage}
\newcommand{\tocback}{\par\vspace{-3pt}{\hfill\footnotesize\hyperlink{toc}{\textcolor{wdgray}{\ensuremath{\uparrow}\,回到目录}}}\par\vspace{-4pt}}

% ---------- 线索栏 ----------
\newcommand{\cueblock}[2]{\Needspace{4\baselineskip}\leavevmode#1{#2}\par\vspace{-\baselineskip}\vspace{-\parskip}}
% 线索直接贴在段首行的左侧（不用 marginnote，避免分页时偶尔落进正文栏）
\newcommand{\cuebox}[1]{\leavevmode\llap{\smash{\begin{minipage}[t]{\marginparwidth}\raggedright\small\linespread{1.18}\selectfont #1\end{minipage}}\hspace{\marginparsep}}}
\newcommand{\cuenote}[1]{\cuebox{{\color{wdblue}\bfseries Q\,}\color{wddark}#1}}
\newcommand{\cuekey}[1]{\cuebox{\color{wdred}\bfseries #1}}

% ---------- 讲解框 ----------
\tcbset{wdbox/.style={enhanced,breakable,arc=2pt,boxrule=0pt,leftrule=3pt,left=7pt,right=7pt,top=5pt,bottom=5pt,
   before skip=7pt,after skip=7pt,fonttitle=\bfseries\small,attach title to upper={\quad},coltitle=#1,
   colframe=#1,colback=#1!6}}
\newtcolorbox{defbox}[1][]{enhanced,breakable,arc=3pt,boxrule=0.8pt,colframe=wdblue,colback=wdblue!4,
   left=8pt,right=8pt,top=6pt,bottom=6pt,fonttitle=\bfseries\small,coltitle=white,colbacktitle=wdblue,
   title={定义#1},before skip=8pt,after skip=8pt}
\newtcolorbox{bridgebox}{wdbox=wdblue,title={承上启下}}
\newtcolorbox{plainbox}{wdbox=wddeepgreen,colframe=wdgreen,colback=wdgreen!8,title={大白话}}
\newtcolorbox{deepbox}{wdbox=wdpurple,title={背后的思想}}
\newtcolorbox{exambox}{wdbox=wdred,title={考点}}
\newtcolorbox{pitfallbox}{wdbox=wdorange,title={易错提醒}}
\newtcolorbox{bookbox}{wdbox=wddark,colback=wdgray!10,colframe=wdgray,title={王道书补充}}
\newtcolorbox{teacherbox}{wdbox=wdorange,colback=wdgold!8,colframe=wdgold,coltitle=wdorange,title={老师原话}}
\newtcolorbox{fixbox}{wdbox=wdred,colback=wdred!4,title={勘误 / 说明}}
% 习题笔记用
\newtcolorbox{questionbox}[1][]{enhanced,breakable,arc=3pt,boxrule=0.8pt,colframe=wddark,colback=white,
   left=8pt,right=8pt,top=6pt,bottom=6pt,fonttitle=\bfseries\small,coltitle=white,colbacktitle=wddark,
   title={#1},before skip=10pt,after skip=6pt}
\newtcolorbox{videobox}{wdbox=wdorange,colback=wdgold!8,colframe=wdgold,coltitle=wdorange,title={视频讲解}}
\newtcolorbox{jiexibox}{wdbox=wddark,colback=wdgray!10,colframe=wdgray,title={书上解析}}
\newtcolorbox{addbox}{wdbox=wdpurple,title={我的补充}}
\newtcolorbox{recallbox}{wdbox=wdblue,title={回顾正课}}
\newtcolorbox{watchbox}{enhanced,breakable,arc=3pt,boxrule=1pt,colframe=wdorange,colback=wdgold!12,
   fonttitle=\bfseries\small,coltitle=white,colbacktitle=wdorange,title={建议去看视频},left=8pt,right=8pt,before skip=8pt,after skip=8pt}
\newtcolorbox{summarybox}{enhanced,arc=3pt,
   boxrule=1pt,colframe=wdblue,colback=white,colbacktitle=wdblue,coltitle=white,fonttitle=\bfseries,
   title={总结区 · 用自己的话复述这一节},left=10pt,right=10pt,top=8pt,bottom=8pt,before skip=14pt}

% ---------- 代码 ----------

\makeatletter
\@ifundefined{Shaded}{}{\renewenvironment{Shaded}{\def\FrameCommand{\fboxsep=5pt\fboxrule=0.5pt\fcolorbox{wdblue!35}{codebg}}%
   \MakeFramed{\advance\hsize-\width\FrameRestore}\small}{\endMakeFramed}}
\makeatother
\fvset{fontsize=\small}

% ---------- 表格 ----------
\renewcommand{\arraystretch}{1.42}
\newcolumntype{Y}{>{\raggedright\arraybackslash}X}

% ---------- 图片 ----------
% 课件截图做成浮动体：LaTeX 会自动用正文填满页面，避免大块空白
\setcounter{topnumber}{3}\setcounter{bottomnumber}{2}\setcounter{totalnumber}{4}
\renewcommand{\topfraction}{0.85}\renewcommand{\bottomfraction}{0.6}\renewcommand{\textfraction}{0.12}\renewcommand{\floatpagefraction}{0.75}
\setlength{\textfloatsep}{10pt plus 3pt minus 2pt}\setlength{\floatsep}{8pt plus 2pt}
\newcommand{\slidefig}[3]{\begin{figure}[!ht]\centering\begin{tcolorbox}[enhanced,width=#2\linewidth,arc=2pt,boxrule=0.6pt,colframe=wdgray!70,
   colback=white,left=2pt,right=2pt,top=2pt,bottom=2pt,drop fuzzy shadow=black!25]\includegraphics[width=\linewidth]{#1}\end{tcolorbox}
   \vspace{-3pt}{\small\color{wdgray}#3}\end{figure}}
\newcommand{\plainfig}[3]{\par\noindent\begin{minipage}{\linewidth}\centering\includegraphics[width=#2\linewidth]{#1}\par\vspace{2pt}{\small\color{wdgray}#3}\end{minipage}\par\vspace{4pt}}

\begin{document}
% ---------- 封面条（仿课件封面） ----------
\thispagestyle{plain}
\begin{tcolorbox}[enhanced,enlarge left by=-4.8cm,width=\linewidth+4.8cm,arc=0pt,boxrule=0pt,
   colback=wddark,interior style={left color=wddark,right color=wdblue},left=18pt,right=18pt,top=14pt,bottom=12pt,
   overlay={\draw[white,line width=2.5pt] ([xshift=10pt,yshift=-10pt]frame.north west) -- ++(0,-0.9cm);
            \draw[white,line width=2.5pt] ([xshift=10pt,yshift=-10pt]frame.north west) -- ++(1.6cm,0);
            \draw[white,line width=2.5pt] ([xshift=-10pt,yshift=10pt]frame.south east) -- ++(0,0.9cm);
            \draw[white,line width=2.5pt] ([xshift=-10pt,yshift=10pt]frame.south east) -- ++(-1.6cm,0);}]
  {\color{white!85!wdlightblue}\small 数据结构 · 第1章
绪论 · 康奈尔听课笔记}\par\vspace{6pt}
  {\color{white}\fontsize{22}{28}\selectfont\bfseries 1.1
课后习题：数据结构的基本概念}\par\vspace{6pt}
  {\color{white!80!wdlightblue}\small 王道书 §1.1.3 试题精选 · §1.1.4
答案与解析 · 习题讲解字幕 P1（选择 1--5）}
\end{tcolorbox}
\vspace{4pt}
\hypertarget{ux8fd9ux7ec4ux9898ux5728ux8003ux4ec0ux4e48}{%
\section{这组题在考什么}\label{ux8fd9ux7ec4ux9898ux5728ux8003ux4ec0ux4e48}}

\Needspace{4\baselineskip}

\cuenote{这 7 道题围绕哪个核心概念？}7
道题几乎都在考同一件事：\hlred{数据结构的三要素}，也就是逻辑结构、存储结构（物理结构）、数据的运算，以及三者之间谁依赖谁。正课
1.1 笔记第二节的这张图，做题时一直放在脑子里就够了：

\begin{figure}[!ht]\centering\adjustimage{scale=0.6,max width=\linewidth,max height=0.6\textheight}{/mnt/project-files/notes/1数据结构/笔记/_build/cache/36f6aaeabd8783752e452c96d248310d3655e8a3.pdf}\end{figure}

{\rowcolors{2}{wdblue!5}{white}\arrayrulecolor{wdblue!35}
\begin{xltabular}{\linewidth}{>{\hsize=1.000\hsize\raggedright\arraybackslash}X>{\hsize=1.000\hsize\raggedright\arraybackslash}X>{\hsize=1.000\hsize\raggedright\arraybackslash}X}
\rowcolor{wdblue}\textcolor{white}{\bfseries 题号} & \textcolor{white}{\bfseries 考点} & \textcolor{white}{\bfseries 答案} \\
\endfirsthead
选择 01 & 三要素有哪些 & C \\ \hline
选择 02 & 线性 / 非线性结构的划分 & A \\ \hline
选择 03 & 一个名字是逻辑结构还是存储结构 & C \\ \hline
选择 04 & 逻辑结构与存储结构的依赖关系 & A \\ \hline
选择 05 & 存储数据时要存什么 & C \\ \hline
综合 01 & 逻辑、物理结构都相同，是不是同一种数据结构 & 不一定 \\ \hline
综合 02 & 同一逻辑结构、不同存储方式，运算效率不同 & 举例 \\ \hline
\end{xltabular}\par\vspace{4pt}}

\begin{fixbox}

\textbf{这节的视频不用看。} 习题讲解视频只讲了选择题
1--5，内容是几个生活类比（容器、学号名单和宿舍），没有正课和书上之外的新知识点，也没有需要截图保存的图。两道综合题视频没讲，下面按书上解析整理，并补了两处说得不严谨的地方。

\end{fixbox}

\hypertarget{ux4e00ux5355ux9879ux9009ux62e9ux9898}{%
\section{一、单项选择题}\label{ux4e00ux5355ux9879ux9009ux62e9ux9898}}

\Needspace{8\baselineskip}\vspace{22pt}

\cueblock{\cuenote}{\relax{三要素是哪三个？缺一个行不行？}}

\begin{questionbox}[01]

一个完整的数据结构通常应包含以下哪些要素（　）？

A. 数据元素及其存储方式\\
B. 数据的逻辑结构和物理结构\\
C. 数据的逻辑结构、存储结构以及在其上定义的基本操作\\
D. 数据对象和数据元素之间的关系

\end{questionbox}

\hlred{\textbf{答案：C}}

\begin{videobox}

老师把数据结构比作\textbf{装数据的容器}：瓶子、盒子这种“形状”是逻辑结构；纸盒还是铁盒、用什么材料做出来是存储结构；倒水、装水是基本操作。存储结构也叫物理结构。

\end{videobox}

\begin{jiexibox}

一个完整的数据结构由三部分组成：逻辑结构（如线性、树形等）、存储结构（物理表示）以及在其上定义的基本操作（如插入、删除等）。只有选项
C 同时包含这三个核心要素。

\end{jiexibox}

\begin{addbox}

B
是最常见的错选，它漏掉了运算。三要素里“数据的运算”和“基本操作”是同一个东西，课件叫运算，题目叫基本操作。综合题
01 会专门考“只看逻辑和物理结构够不够”。

\end{addbox}

\Needspace{8\baselineskip}\vspace{22pt}

\cueblock{\cuenote}{\relax{线性结构的判定标准？集合算线性还是非线性？}}

\begin{questionbox}[02]

下列四种数据结构中，（　）是非线性数据结构。

A. 树　　　　B. 字符串　　　　C. 队列　　　　D. 栈

\end{questionbox}

\hlred{\textbf{答案：A}}

\begin{videobox}

线性结构：除第一个和最后一个元素外，每个元素都\uln{有且仅有一个前驱和一个后继}，是一对一的关系，比如字符串
“hello”。不满足这一条的都是非线性结构，比如集合、树、图。树是一对多：一个领导下面带好几个人。

\end{videobox}

\begin{jiexibox}

树和图是典型的非线性数据结构，其他选项都属于线性数据结构。

\end{jiexibox}

\begin{recallbox}

正课里王道书的分类，选择题最爱考这一张：

\begin{itemize}
\tightlist
\item
  \textbf{线性结构}：一般线性表；受限线性表（\uln{栈、队列、串}）；线性表推广（数组）。
\item
  \textbf{非线性结构}：\hlred{集合}、树形结构、图状结构。
\end{itemize}

栈和队列只是“操作受限”，元素之间仍是一对一，所以还是线性的；串是元素为字符的线性表。

\end{recallbox}

\slidefig{assets/1.1-p15.jpg}{0.55}{课件：四种逻辑结构（正课 1.1）}

\Needspace{8\baselineskip}\vspace{22pt}

\cueblock{\cuenote}{\relax{怎么一眼判断一个名字是逻辑结构还是存储结构？}}

\begin{questionbox}[03]

下列选项中，属于逻辑结构的是（　）。

A. 顺序表　　　　B. 哈希表　　　　C. 有序表　　　　D. 单链表

\end{questionbox}

\hlred{\textbf{答案：C}}

\begin{videobox}

老师用一个班 30 个人举例：按学号 1～30
排成的名单是逻辑结构（而且是有序的线性表，即有序表）；这 30
人实际住在哪间宿舍、坐在教室哪个位置，是物理结构。顺序表、单链表第 2
章学，哈希表在查找那章学，都属于存储结构。

\end{videobox}

\begin{jiexibox}

顺序表、哈希表和单链表是三种不同的数据结构，既描述逻辑结构，又描述存储结构和数据运算。而有序表是指关键字有序的线性表，仅描述元素之间的逻辑关系，它既可链式存储，又可顺序存储，所以属于逻辑结构。

\end{jiexibox}

\begin{addbox}

\textbf{判别方法：看名字里有没有“怎么存”的信息。}

{\rowcolors{2}{wdblue!5}{white}\arrayrulecolor{wdblue!35}
\begin{xltabular}{\linewidth}{>{\hsize=1.000\hsize\raggedright\arraybackslash}X>{\hsize=1.000\hsize\raggedright\arraybackslash}X}
\rowcolor{wdblue}\textcolor{white}{\bfseries 名字里带存储方式 → 存储结构} & \textcolor{white}{\bfseries 名字只说关系 → 逻辑结构} \\
\endfirsthead
顺序表、单链表、双链表、循环链表、静态链表、哈希表（散列表）、顺序栈、链栈、循环队列 & 线性表、有序表、栈、队列、串、树、二叉树、图 \\ \hline
\end{xltabular}\par\vspace{4pt}}

以后遇到“××是一种什么结构”的题，先看它能不能既用数组存、又用指针存。能，就说明它没规定怎么存，是逻辑结构。比如有序表可以用顺序表存，也可以用链表存，所以是逻辑结构。

\end{addbox}

\Needspace{8\baselineskip}\vspace{22pt}

\cueblock{\cuenote}{\relax{谁独立于谁？一种逻辑结构能对应几种存储结构？}}

\begin{questionbox}[04]

下列关于数据结构的说法中，正确的是（　）。

A. 数据的逻辑结构独立于其存储结构\\
B. 数据的存储结构独立于其逻辑结构\\
C. 数据的逻辑结构唯一决定其存储结构\\
D. 数据结构仅由其逻辑结构和存储结构决定

\end{questionbox}

\hlred{\textbf{答案：A}}

\begin{videobox}

还是用班级举例：先有按学号排的名单（逻辑结构），之后排宿舍、排座位都照着这份名单来（存储结构）。所以名单不受宿舍怎么排的影响，A
对；宿舍是照着名单排的，B
错。同一份名单既能排宿舍、又能排座位，存法不止一种，C 错。D
漏了运算，比如可以安排同学上课吃饭，但不能打骂，运算也是有限制的。

\end{videobox}

\begin{jiexibox}

数据的逻辑结构是从面向实际问题的角度出发的，只采用抽象表达方式，独立于存储结构，数据的存储方式有多种不同的选择；而数据的存储结构是逻辑结构在计算机上的映射，它不能独立于逻辑结构而存在。数据结构包括三个要素，缺一不可。

\end{jiexibox}

\begin{addbox}

四个选项正好对应三要素之间的三条关系，记住这三条就能秒选：

\begin{enumerate}
\def\labelenumi{\arabic{enumi}.}
\tightlist
\item
  逻辑结构独立于存储结构；存储结构\hlred{依赖}逻辑结构。（A 对，B 错）
\item
  一种逻辑结构可以有\hlred{多种}存储结构，比如线性表可以存成顺序表，也可以存成链表。（C
  错）
\item
  数据结构要三要素都齐。（D 错，同第 01 题）
\end{enumerate}

\end{addbox}

\Needspace{8\baselineskip}\vspace{22pt}

\cueblock{\cuenote}{\relax{存数据时除了值还要存什么？“存关系”具体怎么存？}}

\begin{questionbox}[05]

在存储数据时，通常不仅要存储各数据元素的值，还要存储（　）。

A. 数据的操作方法　　　　B. 数据元素的类型\\
C. 数据元素之间的关系　　D. 数据的存取方法

\end{questionbox}

\hlred{\textbf{答案：C}}

\begin{videobox}

一个班 30
人可以用一张线性表（数组）存；公司里的上下级关系就得用树来存。关系不同，存法就不同，这正是数据结构这门课的核心。很多人会选
B，觉得 \texttt{int\ a\ =\ 0}
存了类型。但类型在定义里就规定好了，就像数据库的表定义已经写明了每个字段的类型，存每条数据时不用再存一遍。

\end{videobox}

\begin{jiexibox}

在存储数据时，不仅要存储数据元素的值，而且要存储数据元素之间的关系。

\end{jiexibox}

\begin{addbox}

视频对 A、D 只说了句“同样的道理”，这里补全：

\begin{itemize}
\tightlist
\item
  \textbf{A
  操作方法}是算法，也就是程序代码，属于“运算的实现”，不和每个数据一起存。
\item
  \textbf{D 存取方法}（随机存取 /
  顺序存取）是选定存储结构之后\uln{自然产生的性质}，不是额外存进去的东西。顺序存储能随机存取，链式存储只能顺序存取。
\end{itemize}

另外，“存储关系”不一定是多存一份数据。回顾正课的四种存储结构：\textbf{顺序存储}靠物理位置相邻来表示逻辑上相邻，关系藏在位置里；\textbf{链式存储}才是额外存指针，显式地把关系记下来。两种都算存了关系，只是方式不同。

\end{addbox}

\hypertarget{ux4e8cux7efcux5408ux5e94ux7528ux9898}{%
\section{二、综合应用题}\label{ux4e8cux7efcux5408ux5e94ux7528ux9898}}

\Needspace{8\baselineskip}\vspace{22pt}

\cueblock{\cuenote}{\relax{逻辑结构、物理结构都一样，还会是两种不同的数据结构吗？}}

\begin{questionbox}[综合 01]

对于两种不同的数据结构，逻辑结构或物理结构一定不相同吗？

\end{questionbox}

\hlred{\textbf{答案：不一定。}}
两种数据结构的逻辑结构和物理结构可以完全相同，区别在于\uln{运算}不同。

\begin{jiexibox}

应该注意到，数据的运算也是数据结构的一个重要方面。对于两种不同的数据结构，它们的逻辑结构和物理结构完全有可能相同。比如二叉树和二叉排序树，二叉排序树可以采用二叉树的逻辑表示和存储方式，前者通常用于表示层次关系，而后者通常用于排序和查找。虽然它们的运算都有建立树、插入结点、删除结点和查找结点等功能，但对于二叉树和二叉排序树，这些运算的定义是不同的，以查找结点为例，二叉树的平均时间复杂度为
\(O(n)\)，而二叉排序树的平均时间复杂度为 \(O(\log_2 n)\)。

\end{jiexibox}

\begin{fixbox}

书上说二叉排序树查找“平均
\(O(\log_2 n)\)”，前提是树比较平衡。如果按有序序列依次插入，二叉排序树会退化成一条链，查找变成
\(O(n)\)。这在第 7 章查找里会细讲，这里先知道 \(O(\log_2 n)\)
不是无条件成立的。

\end{fixbox}

\begin{addbox}

书上的例子要到第 5、7 章才学，可以先用第 3
章就会学到的\textbf{栈和队列}理解：两者的逻辑结构都是线性结构，都可以用数组顺序存储，但栈只能在一端进出（后进先出），队列一端进、另一端出（先进先出）。逻辑结构和存储结构都一样，只因为运算的规定不同，就是两种数据结构。这也正是第
01 题 B 选项错的原因。

\end{addbox}

\Needspace{8\baselineskip}\vspace{22pt}

\cueblock{\cuenote}{\relax{同一逻辑结构换一种存法，哪些运算变快、哪些变慢？}}

\begin{questionbox}[综合 02]

举例说明对相同的逻辑结构，同一种运算在不同存储方式下实现时，其运算效率不同。

\end{questionbox}

\begin{jiexibox}

线性表既可以用顺序存储方式实现，又可以用链式存储方式实现。在顺序存储方式下，在线性表中插入和删除元素，平均要移动近一半的元素，时间复杂度为
\(O(n)\)；而在链式存储方式下，插入和删除的时间复杂度都是 \(O(1)\)。

\end{jiexibox}

\begin{fixbox}

链表插入、删除是
\(O(1)\)，前提是\uln{已经拿到了要操作位置的前驱结点指针}，此时只改指针。如果题目说的是“在第
\(i\) 个位置插入”，链表得先从头往后找到第 \(i-1\) 个结点，这一步是
\(O(n)\)，整体还是
\(O(n)\)。链表省掉的是\textbf{移动元素}，而不是查找位置。第 2
章算单链表插入删除时会反复用到这一点。

\end{fixbox}

\begin{recallbox}

正课里老师用排队举过一组对照（课件“存储结构的影响”那页），它同时给出了一个反例，作答时两个方向都写上更完整：

{\rowcolors{2}{wdblue!5}{white}\arrayrulecolor{wdblue!35}
\begin{xltabular}{\linewidth}{>{\hsize=1.000\hsize\raggedright\arraybackslash}X>{\hsize=1.000\hsize\raggedright\arraybackslash}X>{\hsize=1.000\hsize\raggedright\arraybackslash}X}
\rowcolor{wdblue}\textcolor{white}{\bfseries 运算} & \textcolor{white}{\bfseries 顺序存储（传统排队）} & \textcolor{white}{\bfseries 链式存储（银行取号）} \\
\endfirsthead
插入、删除 & 后面的元素要整体挪动，\(O(n)\) & 只改指针，\(O(1)\)（已知位置时） \\ \hline
按位置找第 \(i\) 个 & 地址直接算出来，\hlred{随机存取}，\(O(1)\) & 只能从头顺着指针找，\(O(n)\) \\ \hline
\end{xltabular}\par\vspace{4pt}}

\end{recallbox}

\begin{summarybox}

\textbf{这 7 题考的都是三要素。}
逻辑结构独立于计算机，存储结构依赖逻辑结构，而且一种逻辑结构可以有多种存法；数据结构要三要素都齐，逻辑和物理结构都相同的两种数据结构，可以只因运算不同而不同（栈和队列、二叉树和二叉排序树）。判断一个名字是逻辑还是存储结构，看它有没有规定“怎么存”：有序表、栈、队列是逻辑结构，顺序表、链表、哈希表是存储结构。存数据要存值和关系，顺序存储靠位置表示关系，链式存储靠指针。

\textbf{自测}：① 循环队列是逻辑结构还是存储结构？② 链表插入为什么说是
\(O(1)\)，什么情况下不是？③ 字符串、栈、集合，哪个是非线性结构？

\end{summarybox}
\end{document}
